\documentclass[12pt]{article}
\usepackage[T1]{fontenc}
\usepackage[utf8]{inputenc}
\usepackage[brazil]{babel}
\usepackage[margin=1in]{geometry}
\usepackage{ragged2e}
\usepackage[normalem]{ulem}
\setlength{\parindent}{0pt}
\setlength{\parskip}{0.8\baselineskip}
\begin{document}
\justifying
{\Large\bfseries \textbf{PU-TEMA TNU-SOBRESTAMENTO}}\\

\noindent Verifica-se que a questão objeto da controvérsia jurídica suscitada pelo recurso manejado pela parte recorrente encontra-se afetada ao Tema n. 369 da jurisprudência da TNU – PU no processo n. PEDILEF 0001882-94.2021.4.05.8500/SE, pendente de julgamento (afetado em 2024-10-16) –, por meio do qual se elucidará a seguinte questão:

\noindent \textit{"Exclusao de faixa de salario-minimo no calculo de renda per capita para BPC."}

\noindent Nesse contexto, determino o \textbf{SOBRESTAMENTO} do PU interposto nestes autos – o qual versa sobre a mesma matéria –, na forma do\uline{\textbf{ art. 14, II, b}}, da Resolução CJF n. 586/2019 c/c 1.030, III, do CPC, até o julgamento final do representativo da controvérsia, o PU no processo n. PEDILEF 0001882-94.2021.4.05.8500/SE (\textbf{Tema n. }\textbf{369 }\textbf{TNU}).

\noindent Intimem-se. Cumpra-se.
\end{document}
